\begin{afterword}
\summaryhead

The author often meets people who think a superintelligence needs only one simple utility function, such as ``complexity,'' for everything to turn out well. They propose it too fast, and the author's own early term ``supergoal'' may have encouraged them. But a utility function can have many terms, and human values are many: the love of parent for child and of child for parent are not derived from each other or from anything else. Leave out even one value, such as the value of running one's own life, and the result could be a fate worse than death, as in Jack Williamson's ``With Folded Hands.''

Pressed, the proposers say that a superintelligence maximizing complexity would encourage mothers to love their children. That is motivated stopping, since a real complexity maximizer would look for something more complex. It is also a fake morality: defending a principle by its good consequences shows that the consequences are what is valued. Only one's morality reliably reproduces one's moral decisions, so a human morality cannot be compressed into a simple utility function.

\responsehead

The post has two parts: a premise about human values, and an account of how simple value proposals are defended. The second part is the stronger.

Its answer to the proposers is sound. A process that maximized complexity would search for whatever is most complex; it would not stop at a mother's love because that love is complex and people like it. This is the test of the million-dollar laptop in ``Fake Justification,'' applied to a principle instead of a purchase. The post also states correctly what a real argument for one master value would have to show: that other goods are good only when they serve it. That is the same requirement the \textsc{Stanford Encyclopedia of Philosophy} states for utilitarian monists, and some monists, Mill among them, have tried to meet it. The post's target is the weaker proposer who argues from pleasant consequences instead.

The premise is asserted more firmly than it is argued. That human values are many and underived is called a descriptive ``known fact.'' The support is by reference to earlier posts on evolution and desire, chiefly ``Thou Art Godshatter,'' whose central step our notes on that post describe as an inference presented as history. No study of how people's values are related is cited. The corresponding view about what is valuable, value pluralism, is old in moral philosophy; W. D. Ross held it in 1930.

In short: a sound account of how simple utility functions are defended and what a real defence would need, resting on a premise about the many-valued nature of human wants that is called a known fact and argued only in earlier posts.
\end{afterword}
